\subsection{天使大人不在时}

没想到会让亲戚朋友以外的人进屋，周很庆幸平时就有注意保持家里整洁，感谢养成了这个习惯的过去的自己。\\

慧老老实实跟在周他们身后，显得有些坐立不安地跟着周他们从玄关进屋。从把鞋子摆放整齐、走上走廊的动作来看，他似乎受过良好的教导，让人称赞一句父母教导有方也不为过。

当这个念头在脑海中隐约浮现时，周的内心十分复杂，但为了避免被真昼察觉，他始终保持着面无表情。\\

真昼与周不同，她带着一种与其说是冰冷、不如说是失去了温度的表情，默默地走着。

接着，把慧请进客厅，正准备坐下来交谈时，身后传来了一阵「咕——」的、像是地鸣般却又有些可爱的声音。

周心想是怎么回事并回过头，就看到慧一脸尴尬的表情。\\

刚才的声音是什么？\\

是非常熟悉的声音。没错，那就像是有真昼做的料理摆在眼前时，周自然而然会发出的声音——\\

「对不起，请不要在意我的肚子」\\

毕竟那么大声的肚子叫还是挺让人害羞的吧，看着慧慌张地捂着肚子、脸颊微红的样子，周的嘴角浮现出苦笑。

周看向目瞪口呆的真昼身后的时钟，发现已经是平时做饭的时间了。对于正在长身体的孩子来说，这个时间肚子饿也不奇怪，就连周虽然肚子还没叫出声来抗议，但也感到有些饿了。\\

慧的表情仿佛在说「怎么偏偏在这个时候」，因羞耻而微微发抖，但他的肚子却像是不顾及他的心情一般，再次传出了哀怨的叫声。

紧张感瞬间烟消云散，一种难以言喻的尴尬气氛从慧那边飘了过来，但周轻轻一笑，像是要驱散停滞的空气般拍了拍手。\\

「我们吃完饭再谈吧」

「咦？」

「接下来要谈的，对你们两个人来说都是很重要的事情，这个认知没错吧？」

「是、是的」

「那么首先填饱肚子吧。我不认为在能量不足的情况下能集中注意力，肚子饿的话也会失去从容。虽然你应该很在意各种事情，但要想冷静地交谈，最好还是先创造一个双方身心都比较圆满的状态」\\

周觉得，把谈话推迟虽然也会有精神上的负担，但如果在没有余力的状态下听人说话，对后续的影响会更大。

而且，看真昼的样子，他认为还是给她一点时间去做好了解情况的心理准备比较好。\\

「慧，你能接受别人亲手做的饭菜吗？如果不习惯，我就去买点什么」

「没、没关系的……可以吗？」

「我是觉得无所谓，而且放任你饿着肚子不管，我也会良心不安的」\\

慧捂着肚子，脸一下红了，但或许是觉得无法再反驳周的话，他乖乖地保持了沉默。

周对着一直保持安静的真昼温柔地微笑着。\\

「……真昼，你还好吗？如果你觉得无法跟他待在一起，我觉得你要不就先回自己家等着」\\

如果真昼说不愿意，周也考虑让她先回家等候，当作是平复心情的时间——不过，真昼缓缓地摇了摇头。\\

「……没关系。那个，虽然我有很多想法，到现在也有无法接受的地方，但放任一个肚子饿的孩子不管，我心里也不舒服」

「孩、孩子……我已经不是被当成孩子对待的年纪了」

「你几岁了？」

「……13岁」

「那就是孩子啊」

「呜……」\\

在大人看来周也还是个孩子，但他不去考虑这一点，依然把对方当作年纪还小的孩子来对待。

13岁。

大概是初一初二，和真昼相差四岁。

周虽然有很多想法，但这不该由他来说出口。要说的话，也该由真昼来说。\\

「那么，怎么办？如果不介意我们做的饭菜，我们就去准备晚饭了」

「……那就恭敬不如从命了」

「你懂得挺难的词汇啊」

「请、请不要把我当小孩子」\\

慧气鼓鼓地吊起眉角，看上去依然很稚气。

如果把这说出口，他估计就要闹别扭了，于是周轻笑着催促慧到客厅的沙发上坐下。

接着，周轻轻碰了碰真昼的后背，反过来悄悄把她推向玄关的方向。

面对她疑惑的眼神，周尽量用温和的眼神和气氛回应，为了让她暂时离开慧的身边，两人慢慢来到玄关处，周尽量压低声音开口说道。\\

「真昼，你先去换身衣服吧。我来做饭，你慢慢来。等你冷静下来再来我家。好吗？」

「可是」

「现在你的脸色很差，首先从冷静下来开始吧。……其实我是想陪着你的，但从你的角度来看，把他一个人留在这里也不好吧？以真昼的性格，那样反而会静不下心来的。抱歉啊，不能陪着你」\\

老实说，周觉得就这样让真昼一个人待着并不好，但让对慧抱有警戒和怀疑的真昼把慧一个人留在周的家里，她心里大概也会觉得不舒服吧。

综合考虑之后，周得出结论，为了让她暂时冷静下来，还是安排一段独处的时间比较好。

而且，在真昼不在的时候，周也有想问的事情。\\

「那是把我这边的麻烦卷进来的我不好」

「没关系，我完全不觉得……或者说我并没有觉得讨厌。真昼能坦率地依赖我，我很高兴」\\

他甚至觉得这是一种进步，往往把事情藏在心里独自消化的真昼，这次能够坦率地依赖他了。

虽然这也是伴随着巨大冲击和痛苦的事态的证明，所以不能算是一件好事，但她愿意伸出手向周求助这件事，让周感到高兴。\\

「慢慢来就好。首先请优先整理好真昼自己的心情。没关系，有我在你身边……虽然物理上可能会离开几米远？但心是贴在一起的」

「……真是的，周君」\\

不知道是心情稍微好转了，还是被周稍微的玩笑转移了注意力，真昼轻声笑了笑，把头轻轻地靠在了周的手臂上。\\

「我会尽量马上回来的。……并没有到需要周君担心的地步」

「谁知道呢。真昼你可是个爱逞强的人啊。别勉强，好好地把心情平复下来再过来吧」\\

真昼用头咕噜咕噜地向周施压，周也轻声笑了笑，带着鼓励的意味，轻轻握住了真昼的手。\\

送走真昼后，周换了身衣服站在厨房里，而在沙发上待命的慧不知是因为怎么也无法消除紧张，还是觉得不自在，他坐得笔直，却又缩着肩膀。\\

看着他用那种背部像插了尺子一样笔直的漂亮姿势缩成一团，周也觉得很不自在。

在别人家里，而且是在要谈重要的事情之前，再加上还有一个怎么想对自己都没有好感的男人在场，在这样的环境下什么感觉都没有才奇怪，不过周这边是连一丁点要加害他的意思都没有的。

说起来，就算周对他没有好感，但周本身也并不否定他，只是他估计不会这么想。\\

「……嗯，我又不会吃了你。虽然不指望你能完全放松，但你可以随意一点的」

「这、这我知道」\\

觉得他一直保持那个姿势肯定很难受，于是出声搭话后，慧猛地抖了一下身体，然后怯生生地轻轻靠在了沙发的靠背上。

虽然觉得让小孩子这样顾虑挺不好意思的，但从他的立场来说确实没有不紧张的理由，周反而有些后悔自己开口可能让他更加在意了。

不过一直紧张果然还是会累，慧的表情稍微缓和了一些，但依然像是在观察这边一样看着周。\\

我又不会做什么奇怪的事……周一边苦笑，一边把洗好的米放进电饭煲，按下了煮饭开关。\\

「你没有觉得疑惑吗？」\\

伴随着「滴」的一声机械的电子音，周开口问道，慧再次将视线转向这边。\\

「什么？」

「不，就是，我存在本身？你没想过为什么是我在主导局面之类的吗？」

「……如果只有我一个人的话，姐……姐姐，可能连话都不愿意听我讲」

「是啊」

「而且，我给姐姐带来了讨厌的事情，这是事实。她有所警戒，让第三方夹在中间也不奇怪。既然能让姐姐处于愿意听我说话的状态，哥哥你一定是非常亲近的人……是姐姐的男朋友吧？」

「是啊」\\

都已经是初中生了，那也很容易想到，待在（名义上的）姐姐身边表现得很亲密的男人，自然就是男朋友。\\

「而且……那个，母亲的，调查书里，也有你的照片」

「调查书，啊」\\

因为听到了意想不到的词汇，周制作生姜烧腌料的手停了下来。不过仔细想想，真昼的父亲朝阳似乎就在一定程度上调查过周的事情，那么母亲小夜去调查也不奇怪。虽然就个人而言这并不是什么值得高兴的事。\\

「也就是说，你母亲是知道我的存在的，对吧」

「是这样没错」

「在不知道的地方被人擅自调查，感觉可不太好啊。算了，调查女儿周围的情况倒也不是什么奇怪的事，但老实说，知道自己被调查了还是让我吃了一惊」

「吃惊？」\\

确实，调查接近女儿的男人的身份是常有的事，但这大多是出于对女儿的关心才会去做的事。正因为如此，周才对小夜调查这边的事感到惊讶。\\

「嗯，怎么说呢，太意外了吧。我一直以为那个人对真昼完全没有兴趣的」\\

周并没有直接和小夜说过话。只是单方面地听到过她和真昼说话，以及听真昼提起过她而已。

即使是从这种微薄的、甚至算不上关联的仅仅是认知程度的了解上来看，小夜也是一个不会在真昼身上花费时间和精力、不会对她产生兴趣的人。\\

「……从你的角度来看是这样的啊」

「因为没直接说过话所以我也不能说太多，但在我的印象中是这样的。对你来说可能是不想听的事情吧」

「没有。从姐姐的角度来看，这是事实，会这么想也是没办法的。待在姐姐身边的你会这么想也是理所当然的。……把麻烦带给你们，对不起」

「虽然算不上困扰，但是，也很难说这对我和真昼来说究竟是不是好事。毕竟可以肯定的是，我们交往之后，至少被调查过一次」\\

周拿出猪里脊肉放入腌料中腌制，一边回想起他们开始交往的契机——体育节。

体育祭是在去年的六月左右，也就是说是在那之后进行的调查。

到底是为了什么去调查真昼和周的周围情况呢？

为了独自生活的女儿——这种在一般人看来充满爱意又轻松的想象，周是没法全盘接受的。他对真昼的父母并没有信任到这种程度，他们也没有做过能获取周信任的事情。

到底有什么目的？\\

「姑且问一下，对你来说，我认识的小夜女士是你的母亲，这个关系没弄错吧？」

「……详细的说明等你们两个都在的时候我再说。对我来说，她确实是相当于母亲，这一点没错」

「这样啊」\\

从这种说法来看，小夜是生母的可能性很低。

虽然可能性很低，但周也曾怀疑过他是不是亲弟弟或同父异母的弟弟，不过他既不像小夜也不像朝阳，如果是生母的话，应该也没有必要说得这么含糊。\\

那么，他很可能就是真昼曾提到过的，小夜的情人带过来的孩子，这大概是，但真相究竟如何呢。\\

如果假设他是情人的孩子，想到小夜竟然优先对待那个没有血缘关系的孩子，周就感到胃里一阵火热，但他总不能把这还未确定的事，对一个无辜的孩子发泄出来。\\

「我不想说你母亲的坏话，但就我个人的意见来说，因为你母亲对真昼的态度很冷淡，所以对于她真的对真昼有兴趣这件事本身，我就感到很惊讶了」\\

原来她有在在意啊——这种惊讶与困惑，还有一种之前都干什么去了的感觉。

如果这是出于对真昼不利的企图，周该如何保护真昼呢。\\

「她是对母亲来说，有血缘关系的女儿，对吧」

「是啊。但我看到的你母亲，大概和展现在你面前的态度完全不同吧。毕竟对不同的人改变态度也是常有的事」

「……我眼中的母亲，总是很温柔，虽然言辞严厉，但会为了我四处奔走」

「这样啊，如果你心中的母亲是那样，那就是那样吧」\\

啊，这是多么得不到回报啊。

想要被看一眼，哪怕只是一次也好，过去的真昼为此拼命地努力。想象着这些，周嘴里涌出了一股苦涩，但他并没有表现出口，而是咽了下去。\\

「……对姐姐来说，是不一样的对吧」

「从真昼的角度来看确实不同吧。我就不对此做更多的评论了」\\

真昼的过去是属于真昼的。周并没有亲眼见证过的，不能随随便便地说出口。而且，这对慧来说也是会让他受伤的事实，必须避免局外人草率地谈论。周并没有想要伤害慧的意思。

关掉水龙头的流水，周也停止了言语，慧稍微沉默了一下，露出仿佛在烦恼的动作后，抬起了头。\\

「……姐姐现在，那个……是一个人生活，对吧」

「是啊」

「母亲对此也是默许的」

「说是默许，嗯，也算是吧」\\

要说的话甚至到了鼓励的程度了，不过对慧还是先保持沉默吧。

「……很奇怪对吧。我知道。这事很奇怪我还是知道的」\\

刚才，慧说过他13岁了。

到了那个年纪，也能明白这其中的扭曲了。全心全意照顾自己的母亲，却有一个真正的女儿，而且还把她放置不管。

母亲越是照顾自己，就越凸显出她放弃抚养亲生女儿的事实。甚至让人怀疑她作为一个好母亲的面孔究竟是不是真的。

正因为如此，慧才会来到这里吧。\\

「姐姐现在，幸福吗？」

「谁知道呢。要确认她幸不幸福的话，必须要问本人才行。不过我已经决定要让她幸福，也在好好努力，在我的角度看来她看起来很幸福哦。如果这不是我自作多情的话」

「那我就是来妨碍姐姐的幸福的了，对吧」

「也许是吧，我想现在她本人正在心里努力妥协吧」

「对不起」

「你向我道歉我也很为难的。……不过，你带来的问题是那种可以一直视而不见的问题吗？如果我的想象没错的话，这应该迟早会浮出水面的问题吧？」

「……呃」\\

假设慧现在没有来拜访真昼，今后也没有产生疑问，在不知道真昼的情况下度过一生，这种可能性似乎并不大。\\

如今他就是知道了真昼的事，并带着疑问找来了。越是长大，就越会察觉到目前状态的扭曲而想要探究。根据发生时期的不同，可能会对真昼造成巨大的影响。\\

「那样的话，我觉得有我在她身边的时候发生问题，对真昼的心理健康来说比较好。一个人面对是很痛苦的事情。而且接下来就是春假了，比较有空余时间。总比在考试季正当中的时候带来这些问题要好得多吧？」

「……因为我的缘故让她感到难受，对不起」

「所以我说你向我道歉我也很为难啦。我本来想说让你向本人道歉……但是如果对真昼本人说，她大概也不得不说『没关系』，所以从我的角度来看，其实不太希望你向她道歉呢」\\

周有自知之明，知道自己在说非常任性的话。

但是，这毫无疑问是周的真心话。\\

「说实话，你把这件事带过来对真昼来说确实是不好的事，但如果说这说明你本身是个坏人的话，那就错了吧？你终究只是在那样的环境下长大的，你自己并没有做什么错事」\\

虽然是理所当然的事，但事情的起因是真昼的父母小夜和朝阳，暂定为小夜情人的儿子的慧并没有责任。

如果说是慧要求他们忽视真昼的，那他也有责任，但从慧的语气来看，他很可能直到最近才知道真昼的存在。

说到底，从真昼和慧的年龄差来看，真昼实质上被放弃抚养应该是小夜那边的意图。把自己出生前，以及婴儿时期母亲的作为怪罪到慧头上，也有些太残酷了。\\

话虽如此，真昼对慧的存在表现出抗拒感也是理所当然的。\\

「真昼也明白这一点，所以一旦你向她道歉，她就会用理性判断这不是你的错，然后忍耐自己讨厌的心情去选择原谅你」\\

真昼基本上性格温和且理智，所以如果慧真心诚意地道歉，真昼很难选择不原谅。错的是小夜而不是慧，用理性压抑感情的结果，可以预见她会接受道歉。即使内心再怎么沉重郁结，她也会将其吞下。

真昼至今为止到底品尝了多少痛苦，又是如何把这一切全部藏在心里的，周并不清楚。但是，至少在周能看到的范围内，他不想让她去做那种忍受痛苦的事情。\\

「我啊，比起你会更在乎真昼，所以我不想让你做会让真昼感到讨厌的事情。虽然这很残酷呢」

「……不，就像我会优先考虑自己的情况一样，哥哥你优先考虑自己的情况也是理所当然的。而且，从别人看来，起因确实在我身上，所以哥哥和姐姐疏远我也是理所当然的吧」\\

即使被周宣告会轻视他，慧也没有畏缩，而是理所当然地接受了。

感觉到他是个比想象中更坚强的孩子，周大大地眨了眨眼。\\

「有没有人夸你很成熟？」

「……因为我正在努力不让父母蒙羞」

「这样啊」\\

不让父母蒙羞，这说明在慧的眼里，父母看起来很出色。

那是一件非常棒的事情，周也为自己的父母感到骄傲，所以对此有所共鸣。但是，与其说是这个——这也让周理解到，小夜，以及周和真昼都不认识的他的父亲，是真的有在好好参与抚养孩子，这让周感到一丝无奈和难以释怀。

无论真昼怎么渴求、怎么努力都没能得到的东西，他却如此轻易地得到了。而且渴求的对象还是同一个人。

想到这里，周的胸口深处就像被罪恶感和不快感炙烤着一样，感到一阵灼热的钝痛，但他知道慧并没有任何罪过。

所以周只是平静地倾听慧的话语。

慧似乎并没有察觉到周内心的变化，只是有些不自在地将视线落在了膝盖下方。
